As the beam enters the patient, the absorbed dose at a point varies with depth. This variation depends on many conditions: Beam energy, Field size, distance from the source, and Beam collimation system. Thus, calculating the dose in a patient involves considering these parameters and others as they affect depth dose distribution. Hence, the dose calculation system needs to establish percentage depth dose variation along the central axis of the beam.

The Dose increases with depth as we move into depth from the surface up to a maximum value and then decreases with depth. The dose region between the surface and the depth ($d_{max}$) of the maximum dose is called the dose build-up region. This increase in dose deposition is due to the increased range of secondary electrons released in the tissue. The larger the photon energy, the larger the range of secondary electrons, and consequently, the larger the depth of dose maximum ($d_{max}$).
\subsection{Beam Parameters}
\subsubsection{Percentage Depth Dose(PDD)}
The Percentage depth dose (PDD) is defined as the quotient, expressed as a percentage, of the absorbed dose at any depth $d$ to the absorbed dose at a reference depth $d_0$, along the central axis of the beam for a fixed source to surface distance. Mathematically, it can be expressed as:
\begin{equation}PDD(d) = \frac{D(d)}{D(d_0)} \times 100\%\end{equation}
Where $D(d)$ is the absorbed dose at depth $d$ and $D(d_0)$ is the absorbed dose at the reference depth $d_0$. The reference depth $d_0$ is typically chosen as the depth of maximum dose ($d_{max}$) for the given beam energy and field size. 
\begin{figure}
    \centering
    \includegraphics[width=0.8\textwidth]{/Users/souvikdas/Desktop/Lab-Reports/4th Sem/PDD and Profile/Figures/TheoryImages/Factors-of-PDD.jpg}
    \caption{Percentage Depth Dose (PDD) curve showing the variation of absorbed dose with depth for a photon beam. The curve illustrates the build-up region, the depth of maximum dose ($d_{max}$), and the subsequent decrease in dose with increasing depth due to both inverse square law and attenuation.}
\end{figure}
\subsubsection{Surface Dose}
The surface dose is the absorbed dose at the surface of the patient. It is an important parameter in radiation therapy as it can affect the skin and superficial tissues. The surface dose is influenced by factors such as beam energy, field size, and the presence of any bolus material. Generally, for photon beams, the surface dose is lower than the dose at depth due to the build-up effect, where secondary electrons generated by photon interactions deposit their energy at some distance from the point of interaction.
\subsubsection{Practical Range}
In the case of an electron beam, the depth dose at the end remains constant after a certain depth due to the bremsstrahlung production. So, the quantity 'Practical range' is defined as the depth at which is the depth of the point where the tangent to the descending linear portion of the curve (at the point of inflection) intersects the extrapolated background.
\begin{figure}[H]
    \centering
    \includegraphics[width=0.5\textwidth]{/Users/souvikdas/Desktop/Lab-Reports/4th Sem/PDD and Profile/Figures/TheoryImages/PDDofElectron.png}
    \caption{PDD curve for an electron beam showing the practical range, which is the depth at which the dose falls to zero. The curve illustrates the rapid dose fall-off after the practical range due to the limited penetration of electrons in tissue.}
\end{figure}
\subsubsection{Beam Flatness}
The beam flatness for the \textbf{photon beam} is assessed by finding the maximum $D_{max}$ and
minimum $D_{min}$ dose point values on the beam profile within the central 80\% of the beam width and
then using the following relationship. Standard linac specifications generally require that F be less than
3\% when measured in a water phantom at a depth of 10 cm and an SSD of 100 cm for the largest field
size available (usually 40 $\times$ 40 cm$^2$).
\begin{equation}F = \frac{D_{max} - D_{min}}{D_{max} + D_{min}} \times 100\%\end{equation}

Flatness for electron beam along major axes at the depth of dose maximum as the separation between
90\% and 50\% dose points on either side of the beam profile for 10$\times$10 cm$^2$ applicator and maximum applicator size.
The tolerance value is 10 mm.

\subsubsection{Beam Symmetry}
The beam symmetry for the photon is usually determined at $z_{max}$, which represents the most sensitive depth for assessment of this beam uniformity parameter. A typical symmetry specification is that any two dose points on a beam profile, equidistant from the central axis point, are within 2\% of each other. Alternately, areas under the $z_{max}$ beam profile on each side (left and right) of the central axis extending to the 50\% dose level (normalized to 100\% at the central axis point) are determined, and symmetry is then calculated from
\begin{equation}Symmetry = \frac{Area_{left} - Area_{right}}{Area_{left} + Area_{right}} \times 100\%\end{equation}

Symmetry for electron beam along major axes at depth of dose maxima in SSD setup (maximum ratio
of absorbed doses at symmetrical points from central beam axis and more than 1 cm inside the 90%
isodose contour).

\subsubsection{Penumbra}
The photon penumbra is typically defined as the distance between the 80\% and 20\% dose
points on a transverse beam profile measured 10 cm deep in a water phantom. The electron penumbra
is usually defined as the distance between the 80\% and 20\% dose points along a major axis at a given
depth. The IEC defines this depth as one-half of the depth of the 80\% dose on the central axis.

